\section{Câu hỏi nghiên cứu}

Bài báo trả lời 14 research questions (RQs), được nhóm thành 4 clusters:

\subsection{Cluster 1: Character Modeling (RQ1--RQ3)}
\begin{itemize}
    \item \textbf{RQ1}: Làm thế nào mô hình hóa AI Character để duy trì danh tính xuyên suốt?
    \item \textbf{RQ2}: Làm thế nào duy trì personality consistency qua nhiều tương tác?
    \item \textbf{RQ3}: Memory system nên được thiết kế thế nào?
\end{itemize}

\subsection{Cluster 2: Social Intelligence (RQ4--RQ7)}
\begin{itemize}
    \item \textbf{RQ4}: Cơ chế xây dựng relationship giữa Character và user?
    \item \textbf{RQ5}: Emotion nên là output hay internal state?
    \item \textbf{RQ6}: World simulation có cần thiết không?
    \item \textbf{RQ7}: Multi-agent interaction tạo emergent behavior?
\end{itemize}

\subsection{Cluster 3: Technical Foundation (RQ8--RQ9)}
\begin{itemize}
    \item \textbf{RQ8}: Context engineering cho character state?
    \item \textbf{RQ9}: Model independence --- consistency across LLMs?
\end{itemize}

\subsection{Cluster 4: Evaluation \& Long-term (RQ10--RQ14)}
\begin{itemize}
    \item \textbf{RQ10}: Development environment cho character harness?
    \item \textbf{RQ11}: Evaluation methodology?
    \item \textbf{RQ12}: Human user experience?
    \item \textbf{RQ13}: Safety, privacy, user control?
    \item \textbf{RQ14}: Long-term interaction challenges?
\end{itemize}

% =============================================================================
% 4. METHODOLOGY